\documentclass[interno]{gpt}

\titolo{Stato dei lavori e revisione del Piano di qualifica}
\data{2026-11-23}
\orario{18:00}{19:00}
\luogo{Discord}
\redattore{guglielmo}
\verificatore{biasini}
%\assenti{}

\componenti{biasini,cescobolla,corro,guglielmo,mazzurana,spagnoli}
\begin{document}

\section{Ordine del giorno}
\begin{enumerate}
  \item Stato dello studio delle tecnologie.
  \item Allineamento sui requisiti dopo l'incontro con il proponente del 18 novembre 2026.
  \item Stato del riscontro del proponente sui requisiti.
  \item Revisione della prima stesura del \textit{Piano di qualifica}.
\end{enumerate}

\clearpage

\section{Svolgimento}

\subsection{Stato dello studio delle tecnologie}
Il gruppo ha fatto il punto sullo studio delle tecnologie necessarie al progetto.
Lo studio riguarda soprattutto Ruby on Rails, per la parte applicativa, e Bruno, lo strumento indicato da Eggon per descrivere e provare le API.
Con Bruno il gruppo dovrà fornire a Eggon le specifiche delle API per le richieste di ferie e permessi.
Lo studio prosegue nei prossimi giorni.

\subsection{Allineamento sui requisiti}
Il gruppo ha ripreso le indicazioni ricevute da Eggon nell'incontro del 18 novembre 2026, per avere una visione comune prima di aggiornare la lista dei requisiti.
Tra i punti principali:
\begin{itemize}
  \item la conversazione con l'assistente è usa e getta per il dipendente, mentre lato amministrazione le conversazioni vengono conservate per il controllo di qualità;
  \item le richieste di ferie e permessi non si modificano né si ritirano: il dipendente può solo cancellarle e reinserirle;
  \item le fonti sono consultabili da tutti i dipendenti, quindi i requisiti sui destinatari dei \textit{bucket} perdono priorità;
  \item la sostituzione di una fonte con una nuova versione deve conservare lo storico delle versioni usate nelle risposte.
\end{itemize}

\subsection{Stato del riscontro del proponente}
Nell'incontro del 18 novembre 2026 Eggon si era impegnata a dare un riscontro scritto sulla lista dei requisiti entro venerdì 20 novembre 2026.
Al momento della riunione il riscontro non è ancora arrivato.
Il gruppo ha quindi discusso come procedere con l'\textit{Analisi dei requisiti} in base a quando il riscontro arriverà.

\subsection{Revisione della prima stesura del Piano di qualifica}
Il gruppo ha rivisto insieme la prima stesura del \textit{Piano di qualifica}.
La revisione ha riguardato le metriche di processo e di prodotto, le relative soglie e la sezione del cruscotto con i valori dello sprint 1.

\section{Decisioni}
\begin{decisioni}
    \decisione{Prosecuzione dello studio di Ruby on Rails e di Bruno.}
    \decisione{Sollecito a Eggon del riscontro sui requisiti. Se il riscontro arriva prima dell'incontro di mercoledì 25 novembre 2026, il gruppo presenterà in quell'occasione una bozza dell'\textit{Analisi dei requisiti}; se arriva solo mercoledì, la stesura viene posticipata per integrarlo.}
    \decisione{Aggiornamento delle \textit{Norme di progetto} con le decisioni di processo più recenti.}
\end{decisioni}

\section{Azioni da intraprendere}
\begin{azioni}
  \azione[\#141]{Stesura verbale interno della riunione odierna.}{guglielmo}{24/11/2026}
  \azione[\#104]{Prosecuzione dello studio delle tecnologie (Ruby on Rails e Bruno).}{cescobolla}{29/11/2026}
  \azione[\#139]{Aggiornamento delle \textit{Norme di progetto}.}{mazzurana}{29/11/2026}
  \azione{Sollecito a Eggon del riscontro sui requisiti.}{biasini}{23/11/2026}
  \azione[\#140]{Bozza dell'\textit{Analisi dei requisiti}, in base ai tempi del riscontro di Eggon.}{biasini,corro,spagnoli}{25/11/2026}
\end{azioni}

\section{Prossima riunione}
Il prossimo appuntamento è l'incontro settimanale con Eggon di mercoledì 25 novembre 2026 alle 17:00.
La data della prossima riunione interna sarà fissata in base al riscontro del proponente sui requisiti.

\end{document}
